\documentclass[12pt]{article}
\usepackage[utf8]{inputenc}
\usepackage[T1]{fontenc}
\usepackage{hyperref}
\title{Sangir Language Resources for K-12 Curriculum Development}
\author{Compiled by Claude Sonnet 5.5 (claude-sonnet-5-5) for Andr\'as Kornai}
\date{}
\begin{document}
\maketitle

\section{Introduction}
Sangir (Sangihe, Sangi; ISO 639-3 sxn) is an Austronesian language of the Sangihe islands of Indonesia (and the Philippines border region), written in the Latin script. Dialect names reported include Manganitu, Siau, Tabukang, Tagulandang and Tahuna. The searches found only a few reference items and no school curriculum materials.

\section{Primary Educational Resources}
No primers, textbooks or school programmes for Sangir were found in the searches.

\section{Linguistic and Reference Materials}
\begin{itemize}
\item Zacot, Fran\c{c}ois-Robert (2022), \emph{Kamus Dasar Bahasa Sangihe--Bahasa Indonesia}, a basic Sangihe--Indonesian dictionary (Glottolog record): \url{https://glottolog.org/resource/reference/id/669281}
\item Wikipedia, ``Sangir language'': \url{https://en.wikipedia.org/wiki/Sangir_language}
\item Wikipedia, ``Sangir people'': \url{https://en.wikipedia.org/wiki/Sangir_people}
\item Austronesian Basic Vocabulary Database, Sangir (Tabukang dialect) word list: \url{https://abvd.eva.mpg.de/austronesian/language.php?id=446}
\item Atlantis Press paper concerning the Tagulandang dialect (content not read): \url{https://atlantis-press.com/article/25903921.pdf}
\end{itemize}
No full grammar was found.

\section{Religious and Community Resources}
None specific to Sangir were found.

\section{Digital and Audio Resources}
\begin{itemize}
\item Verbix, Sangir verb page (short verb list with glosses): \url{https://docs.verbix.com/Languages/Sangir}
\end{itemize}
No audio or video resources were found.

\section{Curriculum Development Context}
Sangir has a large speaker base but is not known to be taught in schools. The Zacot dictionary and the ABVD word list could seed vocabulary, but a curriculum would require new textbooks, a standardized orthography, and local teacher training, in cooperation with Sangihe communities. The listed items were not examined in full.

\end{document}